This section fixes the evaluation scope used throughout the paper. Unless stated otherwise, results use WOMD \texttt{validation\_interactive} (12{,}800 scenarios), $K{=}6$, and report compliance on the \emph{selected} trajectory. Protocol~A is the operating setting for selector comparisons because it matches the proxy stack used by RECTOR; Protocol~B is the completeness reference because it evaluates the full 28-rule catalog with oracle applicability. We distinguish candidate-set quality from selected-trajectory quality because RECTOR changes only the selection step.
%-------------------------------------------------------------------------------
\subsection{Dataset}
\label{subsec:dataset}
%-------------------------------------------------------------------------------

\subsubsection{Waymo Open Motion Dataset (WOMD)}

We use the Waymo Open Motion Dataset (WOMD) v1.3.0~\cite{WOMD_20201}, which provides temporally aligned trajectories for the ego vehicle and surrounding agents, HD map context (lane topology, crosswalks, boundaries), and traffic control state annotations (e.g., signals and stop signs) that are required for both motion prediction and rule applicability/violation computation. Table~\ref{tab:womd_stats} summarizes the key statistics.

\begin{table}[t]
\centering
\caption{Waymo Open Motion Dataset Statistics (Interactive Scenarios)}
\label{tab:womd_stats}
\begin{tabular}{ll}
\toprule
\textbf{Property} & \textbf{Value} \\
\midrule
\multicolumn{2}{l}{\textit{Dataset Overview}} \\
Total augmented scenarios & 160,184 \\
Training split & 102,040 augmented scenarios \\
Validation split & 12,800 evaluation scenarios \\
Test split & 14,665 augmented scenarios (withheld) \\
Interactive criterion & Exactly 2 objects\_of\_interest \\
Geographic regions & Phoenix, San Francisco, Mountain View \\
Traffic convention & Right-hand traffic (US) \\
Average agents per scenario & 8.2 \\
Maximum agents & 40+ \\
Object types & Vehicles, pedestrians, cyclists \\
Agent tracking consistency & Validity masks over the 9.1\,s window \\
\midrule
\multicolumn{2}{l}{\textit{Map Context}} \\
Map coverage & Full HD maps \\
Lane annotations & Complete lane topology \\
Traffic signals & Per-timestep state annotations \\
Stop signs & Location and orientation \\
Crosswalks & Polygon boundaries \\
Speed limits & Where available \\
\bottomrule
\end{tabular}
\end{table}

\subsubsection{Data Split Strategy}

We follow WOMD's official interactive splits to avoid leakage across training and evaluation distributions. Training uses 102,040 augmented interactive scenarios, while the canonical selector benchmark uses the \texttt{validation\_interactive} split (12{,}800 scenarios sampled from 150 validation shards). Each augmented TFRecord entry wraps the original WOMD Scenario proto together with per-rule applicability labels, violation flags, and severity scores computed by the \texttt{waymo\_rule\_eval} pipeline; these labels serve as oracle applicability ground truth in ablations (Table~\ref{tab:applicability_ablation}). The test split is withheld and is not accessed during development.

\subsubsection{Interactive Scenario Selection}

WOMD defines an \emph{interactive} scenario by the presence of exactly two \texttt{objects\_of\_interest}, which serve as the primary interacting entities. This restriction emphasizes cases where inter-agent coupling and priority reasoning are salient, making them an appropriate stress test for rule-aware selection. We use all interaction types present in the validation split: vehicle-to-vehicle~(v2v), vehicle-to-pedestrian~(v2p), vehicle-to-cyclist~(v2c), and other agent pairings.

\subsubsection{Scenario Filtering}

We apply minimal filtering to preserve the natural distribution while removing degenerate cases that do not meaningfully test interaction or map reasoning: scenarios with fewer than 2~agents are excluded, scenarios with ego speed consistently below $0.1\,\text{m/s}$ are excluded, and scenarios with incomplete map features required by the rule set are excluded. All other scenarios are retained to maintain geographic and behavioral diversity. Overall, approximately 98\% of scenarios survive these filters. The Legal-tier rule catalog reflects traffic norms applicable to U.S. right-hand traffic scenarios in WOMD; extending to other jurisdictions would require re-parameterizing the catalog and re-evaluating proxy coverage accordingly.

%-------------------------------------------------------------------------------
\subsection{Evaluation Protocols}
\label{subsec:eval_protocols}
%-------------------------------------------------------------------------------

Compliance numbers are only meaningful when the evaluation scope is unambiguous. We use Protocols~A and~B as defined in Section~\ref{Sec:Rule_Evaluation_And_Data_Pipeline} (Table~\ref{tab:training_eval_comparison}), but they play different roles: Protocol~A is the operating selector benchmark, while Protocol~B is the reference evaluator for full-rule coverage and for diagnosing what the proxy stack suppresses.

\noindent\textbf{Aggregation conventions.}  Three summary statistics are used throughout the results: \emph{Total~Viol.\%} is the fraction of scenarios where the selected trajectory violates at least one rule across the rulebook (union probability); \emph{Per-tier~Viol.\%} is the fraction violating at least one rule in that tier; and \emph{S+L~Viol.\%} is the fraction violating at least one Safety or Legal tier rule. Per-tier rates may sum to more than Total~Viol.\% because a single scenario can violate rules in multiple tiers.

\noindent Protocol~B applies the full rule set with oracle applicability, so its absolute violation rates are higher than Protocol~A. When the paper discusses "full compliance" or "complete rule coverage," Protocol~B is the relevant reference; for deployment-facing interpretation, Protocol~B together with conservative always-on Safety/Legal activation is the more conservative reading of the system than Protocol~A with the learned applicability head.

\noindent\textbf{Evaluation populations.} Several scenario subsets appear in different analyses. Table~\ref{tab:eval_populations} gives one reference for which subset is used where and whether it is part of the main selector benchmark or a diagnostic/ablation study.

\begin{table*}
\centering
\caption{Evaluation Population Reference}
\label{tab:eval_populations}
\setlength{\tabcolsep}{3pt}
\begin{tabular}{lrl}
\toprule
\textbf{Population} & \textbf{$N$} & \textbf{Used in} \\
\midrule
Canonical validation & 12{,}800 & Main selector benchmark; Protocol~A comparisons; headline compliance results \\
Augmented validation & 43{,}219 & Weight-grid, applicability, and cross-evaluation diagnostics; Protocol~B selector comparisons \\
Proxy alignment & 1{,}000 & Proxy-vs-full-evaluator ranking correlation and safety audit \\
Perturbation & 500 & Robustness to perception and map noise \\
Constrained baselines & 1{,}000 & Simplified mode-diversity comparison \\
Waymax bridge & 55 & Log-replay bridge validation; integration sanity check \\
\bottomrule
\end{tabular}
\end{table*}

%-------------------------------------------------------------------------------
\subsection{Evaluation Metrics}
\label{subsec:metrics}
%-------------------------------------------------------------------------------

We report three metric groups: trajectory accuracy, rule compliance, and computational efficiency.

\subsubsection{Trajectory Accuracy Metrics}

We use the standard geometric metrics (ADE, FDE, minADE, minFDE, MissRate) as defined in Section~\ref{subsec:geometric_metrics}, evaluated over the WOMD prediction horizon ($T{=}50$ at 10\,Hz) with miss threshold~$\delta = 2.0\,\text{m}$.  In addition, we report \emph{selected-trajectory} metrics that quantify the accuracy of the final decision after lexicographic selection:
\begin{equation}
\text{selADE} = \text{ADE}(\hat{\tau}, \tau^*), \qquad
\text{selFDE} = \text{FDE}(\hat{\tau}, \tau^*),
\end{equation}
where $\hat{\tau}$ is the trajectory chosen by the selection strategy.  These quantify decision quality, not only candidate-set diversity. 
\subsubsection{Rule Compliance Metrics}

Compliance is evaluated using the pipeline described in Sections~\ref{subsec:rule_catalog}--\ref{subsec:proxy_functions}. We report tier-wise rates for the main paper and per-rule rates for diagnosis.

\textbf{Tier-wise violation rate.}
For tier~$\ell$ with aggregated score~$S_\ell(\tau; s)$:
\begin{equation}
\text{ViolationRate}_\ell
= \frac{1}{N} \sum_{i=1}^{N} \mathbb{1}\!\left[S_\ell(\tau_i; s_i) > 0\right].
\end{equation}

\textbf{High-priority (S+L) violations.}
We treat violations in the top two tiers (Safety~$\cup$~Legal) as high-priority:
\begin{align}
\text{S{+}L\_ViolationRate}
&=\mathbb{E}_{s \sim \mathcal{D}} \left[
\mathbf{1}\!\left[S_0(\tau; s) > 0 \ \text{or}\ S_1(\tau; s) > 0\right]
\right].\nonumber
\end{align}

\textbf{Per-rule violation rate.}
For each rule~$r$ with violation score~$V_r(\tau; s)$:
\begin{equation}
\text{ViolationRate}_r
= \frac{1}{N} \sum_{i=1}^{N} \mathbb{1}\!\left[V_r(\tau_i; s_i) > 0\right].
\end{equation}
\noindent All tables report violation rates for conciseness; the \emph{Compliance Rate} at tier~$\ell$ is the complement, $\text{CR}_\ell = 1 - \text{ViolationRate}_\ell$, representing the fraction of scenarios where the selected trajectory satisfies all rules in that tier.

\subsubsection{Closed-Loop Simulation Metrics}
\label{subsubsec:closed_loop_metrics}

For the Waymax-based closed-loop bridge validation (Section~\ref{subsec:closed_loop}), we collect Waymax's built-in metrics and custom motion-planning metrics. \textbf{Overlap rate} is the fraction of simulation steps where the ego vehicle's bounding box overlaps any other agent. \textbf{Log divergence}~(m) measures displacement between the simulated and logged trajectory. \textbf{Kinematic infeasibility rate} is the fraction of steps whose action violates the dynamics model's feasibility bounds. \textbf{Longitudinal jerk}~(m/s$^3$) quantifies ride comfort via the time derivative of acceleration. \textbf{Minimum clearance}~(m) is the closest Euclidean distance between the ego and any other valid agent at each step. \textbf{Time-to-collision}~(s) estimates the time before the first collision under a constant-velocity assumption. These metrics are included because the bridge-validation subsection reports them.

\subsubsection{Computational Efficiency Metrics}

We report latency, throughput, and memory footprint under the same batching and hardware conditions used for all models.  Inference latency is defined as
\begin{equation}
\text{Latency}_{\text{mean}} = \frac{1}{N} \sum_{i=1}^{N} t_{\text{forward}}^{(i)},
\end{equation}
where $t_{\text{forward}}^{(i)}$ includes the forward pass and the selection step, measured with GPU synchronization; we also report P95~latency. Throughput is the reciprocal:
\begin{equation}
\text{Throughput} = \frac{N}{\sum_{i=1}^{N} t_{\text{forward}}^{(i)}}.
\end{equation}
We additionally report parameter counts and peak GPU memory usage during inference.

%-------------------------------------------------------------------------------
\subsection{Baseline Methods}
\label{subsec:baseline_methods}
%-------------------------------------------------------------------------------

Comparisons serve two purposes. External baselines provide context for candidate-set accuracy relative to prior WOMD predictors, while internal baselines isolate the selector to a fixed candidate set.

\subsubsection{External Trajectory Accuracy Baselines}

We compare RECTOR against five representative motion prediction methods.  These numbers are included as geometric-accuracy context rather than as the paper's causal comparison, which is the internal same-candidate-set study below.  Unless explicitly stated as reproduced in our pipeline, the external baseline values follow the published descriptions or reference reports of the respective methods and may therefore reflect different implementation details even when they target the same \texttt{validation\_interactive} setting.

\textbf{LaneGCN}~\cite{liang2020lanegcn}: graph neural network over lane-graph structure.  Reported metrics: minADE~= 0.87\,m, minFDE~= 1.92\,m, Miss Rate~= 28.4\%, throughput~= 180~samples/s.  \textbf{DenseTNT}~\cite{gu2021densetnt}: dense goal-based prediction with iterative refinement and WTA-style training.  Reported metrics: minADE~= 0.73\,m, minFDE~= 1.55\,m, Miss Rate~= 20.1\%, throughput~= 120~samples/s.

\textbf{M2I}~\cite{Sun_2022}: interaction-centric scene encoding with multi-modal decoding.  Reported metrics: minADE~= 0.75\,m, minFDE~= 1.63\,m, Miss Rate~= 22.8\%.

\textbf{Wayformer}~\cite{nayakanti2023wayformer}: Transformer-based motion forecasting with long-range attention.  Reported metrics: minADE~= 0.68\,m, minFDE~= 1.42\,m, Miss Rate~= 18.9\%, throughput~= 85~samples/s.

\textbf{MTR}~\cite{shi2022motion}: Transformer-based motion modeling with structured candidate generation.  Reported metrics: minADE~= 0.70\,m, minFDE~= 1.48\,m, Miss Rate~= 19.5\%, throughput~= 95~samples/s.

\textbf{RECTOR (Ours):}  Transformer--CVAE with rule-aware lexicographic selection, trained via the two-stage procedure in Section~\ref{subsec:system_overview}. Reported metrics: minADE~= 0.684\,m, minFDE~= 1.270\,m, Miss Rate~= 18.56\%, throughput~= 136.5~samples/s.

\subsubsection{Internal Selection Strategy Baselines}

The more important comparison is the \emph{internal} one: all three strategies below operate on the same $K{=}6$ candidate set, so any difference in compliance is attributable entirely to the selection mechanism.

\textbf{Confidence Only:}  $\hat{\tau} = \arg\max_k p^{(k)}$.  No rule awareness.

\textbf{Weighted Sum:}  $\hat{\tau} = \arg\min_k \sum_r w_r \hat{a}_r \bar{V}_{r,k}$. Uses the same predicted applicability as RECTOR to ensure a fair comparison; per-rule weights are uniform ($w_r{=}1$ for all $r$). On this evaluation, weighted-sum achieves near-identical compliance to RECTOR; the qualitative distinction between the two selectors is summarized in Section~\ref{subsec:strategy_comparison}.

\textbf{Human GT (reference):}  Ground-truth trajectories evaluated under the identical Protocol~A pipeline.  This serves as the \emph{imitation ceiling}---the performance of a system that perfectly replicates human drivers---and its non-zero Legal-tier violation rate~(12.8\%) establishes that human-like trajectories are not synonymous with rule-compliant ones.

%-------------------------------------------------------------------------------
\subsection{Implementation Details}
\label{subsec:implementation}
%-------------------------------------------------------------------------------

To enable reproducibility, we report the training and inference configuration below.

\subsubsection{Model Architecture}

RECTOR is implemented as an interaction-aware encoder paired with a multimodal CVAE trajectory decoder, augmented by a rule-applicability head and a tiered scorer. Table~\ref{tab:param_breakdown} summarizes the configuration and parameter allocation; the additional capacity beyond the motion backbone is concentrated in rule applicability prediction, while the tiered scorer itself is lightweight and primarily structural.
\subsubsection{Training Configuration}

Training uses AdamW with a OneCycleLR cosine schedule and a brief warm-up. Table~\ref{tab:training_config} reports the hyperparameters, loss weights, augmentation settings, and hardware.

\begin{table}[t]
\centering
\caption{Training hyper-parameters and Configuration}
\label{tab:training_config}
\begin{tabular}{ll}
\toprule
\textbf{hyper-parameter} & \textbf{Value} \\
\midrule
\multicolumn{2}{l}{\textit{Optimizer Settings}} \\
Optimizer & AdamW \\
Max learning rate & $3 \times 10^{-4}$ \\
Learning rate schedule & OneCycleLR (cosine) \\
Weight decay & $1 \times 10^{-4}$ \\
Gradient clipping & max norm = 1.0 \\
Batch size & 256 \\
Epochs (Stage~1 / Stage~2) & 20 / 20+5 \\
Precision & Mixed (AMP) \\
Random seed & 42 \\
\midrule
\multicolumn{2}{l}{\textit{Loss Weights}} \\
$\lambda_{\text{WTA}}$ (reconstruction) & 20.0 \\
$\lambda_{\text{end}}$ (endpoint) & 5.0 \\
$\lambda_{\text{KL}}$ (KL divergence) & 0.05 \\
$\lambda_{\text{goal}}$ (goal prediction) & 2.0 \\
$\lambda_{\text{smooth}}$ (smoothness) & 1.0 \\
$\lambda_{\text{conf}}$ (confidence) & 1.0 \\
$\lambda_{\text{app}}$ (applicability) & 1.0 \\
$\lambda_{\text{rule}}$ (proxy compliance) & 1.0 \\
\midrule
\multicolumn{2}{l}{\textit{Checkpoint Selection}} \\

Selection criterion & Best validation loss \\
Early stopping patience & 15 epochs \\
\midrule
\multicolumn{2}{l}{\textit{Data Augmentation}} \\
Random rotation & $\pm 15^\circ$ \\
Position noise & $\mathcal{N}(0, 0.1^2)$\,m \\
Agent dropout & 10\% \\
Time reversal$^*$ & 50\% probability \\
Mirror flip & 50\% probability \\
\midrule
\multicolumn{2}{l}{\textit{Hardware}} \\
GPU & NVIDIA RTX A3000 (12\,GB) \\
Training duration & $\sim$12 hours \\
Peak GPU memory & 9\,GB \\
\bottomrule
\end{tabular}
\vspace{1mm}
\caption*{\footnotesize $^*$Time reversal is applied to spatial position sequences only, not to rule violation labels or temporal rule semantics. Rule applicability and violation labels are computed on the original (non-reversed) trajectory.  Training follows a two-stage procedure: Stage~1 trains only the applicability head (focal BCE, $\gamma{=}2.0$) on frozen encoder features for 20 epochs; Stage~2 fine-tunes the full model end-to-end with all eight loss terms for 20+5 epochs.  The loss weights in this table apply to Stage~2; Stage~1 uses $\mathcal{L}_{\text{app}}$ only.}
\end{table}

\subsubsection{Loss Functions}

The full training objective is a flat weighted sum of eight terms as specified in Section~\ref{subsec:training_objective} ($\mathcal{L}_{\text{WTA}}$, $\mathcal{L}_{\text{end}}$, $\mathcal{L}_{\text{KL}}$, $\mathcal{L}_{\text{goal}}$, $\mathcal{L}_{\text{smooth}}$, $\mathcal{L}_{\text{conf}}$, $\mathcal{L}_{\text{app}}$, $\mathcal{L}_{\text{rule}}$). The checkpoint with the lowest composite validation loss is selected; because the reconstruction term ($\lambda_{\text{WTA}}{=}20.0$) dominates the total, the best validation loss closely tracks the best validation ADE.

\subsubsection{Inference Configuration}

Inference runs in FP32 with batch size~128; selection uses the tiered lexicographic procedure defined in Section~\ref{subsec:lexicographic_selection}. Latency, throughput, and memory footprint are reported in Table~\ref{tab:efficiency} (Section~\ref{Sec:Accuracy_SectionRule-Compliance_Behavior_and_Efficiency}).

%-------------------------------------------------------------------------------
\subsection{Reproducibility}
\label{subsec:reproducibility}
%-------------------------------------------------------------------------------

Beyond the reference training configuration, full reproducibility requires specifying the software environment, random-seed management, and the data-preprocessing pipeline.

\subsubsection{Software Environment}

All experiments use Python~3.10, PyTorch~2.2, CUDA~12.1, and Waymo Open Dataset~v1.5.2; full dependency versions are provided with the code release. The closed-loop bridge validation (Section~\ref{subsec:closed_loop}) additionally uses Waymax~0.1.0~\cite{gulino2024waymax} and JAX~0.6.2 with the CUDA~12 plugin.

\subsubsection{Random Seed Management}

All experiments use a fixed random seed of 42 across PyTorch, NumPy, and Python's \texttt{random} module.  CUDA deterministic mode and cuDNN deterministic mode are both enabled, with an expected $\sim$2\% throughput cost from the latter.

\subsubsection{Data Preprocessing}

The preprocessing pipeline consists of seven stages: parsing TFRecord scenarios from WOMD; extracting ego and agent trajectories with validity masks and dropping invalid timesteps; transforming all states to an ego-centric frame at $t{=}0$; sampling lane polylines within a 100\,m radius and constructing lane-graph features; extracting traffic control states and geometry required by the rulebook; computing ground-truth rule supervision signals from map context and annotated states; and caching processed features as NumPy arrays for fast, deterministic loading.

\subsubsection{Expected Runtime}

The two-stage training requires approximately 12 hours total on a single NVIDIA RTX~A3000 Laptop GPU (12\,GB); full evaluation with rule violation computation on 12{,}800 scenarios takes approximately 9 minutes (536.7\,s); and inference latency is 7.3\,ms per sample (mean) or 9.2\,ms (P95). This modest single-GPU setup enhances reproducibility but also constrains the scale of additional ablations and large-sweep experiments in the current study.
